% =====================================================================
% Chapter 2: Requirements, Technology and Architecture
% =====================================================================
\chapter{Requirements and System Architecture}
\label{ch:requirements}

\section{Functional Modules}
\label{sec:modules}
The application is divided into five functional modules. A module represents a piece of
functionality, not a single table; most modules use several tables
(Table~\ref{tab:modules}).

\begin{table}[H]
  \centering
  \caption{Functional modules and the tables each one uses}
  \label{tab:modules}
  \begin{tabularx}{\textwidth}{@{}lXl@{}}
    \toprule
    \textbf{Module} & \textbf{Purpose} & \textbf{Tables used} \\
    \midrule
    1. Startup Management     & Register, view, search, update and delete startups and their founders
                              & \code{startup}, \code{founder} \\
    2. Application Management & Submit applications, track and change their status, admit
                                accepted startups into incubation
                              & \code{application}, \code{evaluation}, \code{incubation} \\
    3. Evaluation Management  & Maintain evaluators, record scores and feedback, view results
                              & \code{evaluator}, \code{evaluation}, \code{application} \\
    4. Mentor Management      & Maintain mentors, assign them to incubated startups, maintain
                                mentorship records
                              & \code{mentor}, \code{mentorship}, \code{incubation} \\
    5. Progress Management    & Create milestones, update their status, track progress
                              & \code{milestone}, \code{incubation}, \code{mentorship} \\
    \bottomrule
  \end{tabularx}
\end{table}

\section{Functional Requirements}
\label{sec:fr}
Each requirement has an identifier (FR-$x.y$) so that it can be traced to tables, queries and
test cases in later chapters.

\begin{longtable}{@{}lp{0.80\textwidth}@{}}
  \caption{Functional requirements}\label{tab:fr}\\
  \toprule
  \textbf{ID} & \textbf{Requirement} \\
  \midrule
  \endfirsthead
  \toprule
  \textbf{ID} & \textbf{Requirement} \\
  \midrule
  \endhead
  \multicolumn{2}{@{}l}{\textit{Module 1 --- Startup Management}} \\
  FR-1.1 & Register a new startup. \\
  FR-1.2 & Add one or more founders to a startup. \\
  FR-1.3 & View the list of startups, with search by name or city and filter by sector. \\
  FR-1.4 & View the full details of one startup, including its founders. \\
  FR-1.5 & Update startup and founder information. \\
  FR-1.6 & Delete a startup only where appropriate (see BR-9). \\
  \midrule
  \multicolumn{2}{@{}l}{\textit{Module 2 --- Application Management}} \\
  FR-2.1 & Submit an incubation application for a registered startup. \\
  FR-2.2 & View all applications, with filter by status. \\
  FR-2.3 & Track the status of a given application. \\
  FR-2.4 & Update application status (Submitted $\rightarrow$ Under Review $\rightarrow$ Accepted / Rejected). \\
  FR-2.5 & On acceptance, admit the startup into incubation as a single transaction. \\
  \midrule
  \multicolumn{2}{@{}l}{\textit{Module 3 --- Evaluation Management}} \\
  FR-3.1 & Add, view, update and delete evaluators. \\
  FR-3.2 & Record an evaluation score for an application by an evaluator. \\
  FR-3.3 & Store written feedback with each evaluation. \\
  FR-3.4 & View evaluation results per application (each score, average, highest, lowest). \\
  \midrule
  \multicolumn{2}{@{}l}{\textit{Module 4 --- Mentor Management}} \\
  FR-4.1 & Add, view, update and delete mentors. \\
  FR-4.2 & Assign a mentor to an admitted (incubated) startup. \\
  FR-4.3 & View mentor assignments by mentor and by startup. \\
  FR-4.4 & Maintain mentorship records (assignment date, end date, status, remarks). \\
  \midrule
  \multicolumn{2}{@{}l}{\textit{Module 5 --- Progress Management}} \\
  FR-5.1 & Create milestones for an incubated startup. \\
  FR-5.2 & Update milestone status (Pending $\rightarrow$ In Progress $\rightarrow$ Completed). \\
  FR-5.3 & Record the completion date when a milestone is completed. \\
  FR-5.4 & View the progress of a startup (completed versus total milestones, percentage). \\
  FR-5.5 & Identify overdue milestones. \\
  \bottomrule
\end{longtable}

\section{Business Rules}
\label{sec:br}
Business rules are the real-world rules that the stored data must obey. During design, each
rule becomes a constraint, a relationship cardinality or a transaction
(Table~\ref{tab:br}). How each rule is enforced is explained in Chapter~3.

\begin{table}[H]
  \centering
  \caption{Business rules}
  \label{tab:br}
  \begin{tabularx}{\textwidth}{@{}lX@{}}
    \toprule
    \textbf{ID} & \textbf{Rule} \\
    \midrule
    BR-1  & A startup has one or more founders; a founder belongs to exactly one startup. \\
    BR-2  & Startup names and e-mail addresses are unique. \\
    BR-3  & A startup may apply several times (e.g.\ after a rejection) but may have only one
            open (Submitted or Under Review) application at a time. \\
    BR-4  & Application status is one of Submitted, Under Review, Accepted or Rejected. \\
    BR-5  & An application may be evaluated by many evaluators, but each evaluator evaluates a
            given application at most once. \\
    BR-6  & An evaluation score is a whole number from 0 to 100. \\
    BR-7  & An application can be accepted only if it has at least one evaluation; accepting it
            creates an incubation record, and both changes succeed or fail together. \\
    BR-8  & Mentors can be assigned only to incubated startups; a mentor can guide many startups
            and a startup can have many mentors. \\
    BR-9  & A startup with application history cannot be deleted; a startup without applications
            can be deleted together with its founders. \\
    BR-10 & A milestone is Pending, In Progress or Completed; a completion date exists if and only
            if the status is Completed. \\
    BR-11 & Dates must be logical (for example, the expected end date is after the start date). \\
    \bottomrule
  \end{tabularx}
\end{table}

\section{Reports and Outputs}
\label{sec:reports}
The application produces the reports in Table~\ref{tab:reports}. Each report also exercises a
particular SQL concept from the course.

\begin{table}[H]
  \centering
  \caption{Reports produced by the application}
  \label{tab:reports}
  \begin{tabularx}{\textwidth}{@{}lXl@{}}
    \toprule
    \textbf{ID} & \textbf{Report} & \textbf{Main SQL concept} \\
    \midrule
    R1  & Startups per sector                                  & \code{COUNT}, \code{GROUP BY} \\
    R2  & Applications and funding requested by status         & \code{COUNT}, \code{SUM}, \code{AVG} \\
    R3  & Evaluation summary per application                   & \code{AVG}, \code{MIN}, \code{MAX}, joins \\
    R4  & Applications scoring above the overall average       & subquery in \code{HAVING} \\
    R5  & Applications not yet evaluated                       & \code{LEFT JOIN \ldots\ IS NULL} \\
    R6  & Mentor workload (active startups per mentor)         & \code{LEFT JOIN}, \code{COUNT} \\
    R7  & Incubated startups without an active mentor          & \code{LEFT JOIN} \\
    R8  & Progress per startup (\% milestones completed)       & \code{CASE} inside \code{SUM} \\
    R9  & Overdue milestones                                   & date arithmetic \\
    R10 & Highest-scoring evaluation of each application       & correlated subquery \\
    R11 & Incubation dashboard                                 & view \\
    \bottomrule
  \end{tabularx}
\end{table}

\section{Non-Functional Requirements}
\begin{itemize}
  \item The application runs locally on a single Windows computer.
  \item Data integrity is enforced in the database, not only in the user interface.
  \item All SQL is written by hand (no ORM) and every user-supplied value is passed as a query
        parameter, which prevents SQL injection.
  \item The code is kept simple and readable, with one file per functional module.
\end{itemize}

\section{Technology Stack}
\label{sec:tech}
\begin{table}[H]
  \centering
  \caption{Technology stack}
  \label{tab:tech}
  \begin{tabularx}{\textwidth}{@{}llX@{}}
    \toprule
    \textbf{Component} & \textbf{Version} & \textbf{Role in the project} \\
    \midrule
    PostgreSQL   & 18.6   & Relational database server: tables, constraints, views, transactions \\
    Python       & 3.14   & Application logic \\
    psycopg      & 3.3.8  & Database driver: sends SQL and parameters to PostgreSQL and returns rows \\
    Streamlit    & 1.67.0 & Web user interface: forms, tables, charts and tabs in the browser \\
    psql, pgAdmin 4 & ---  & Running SQL scripts and inspecting the database \\
    \bottomrule
  \end{tabularx}
\end{table}

\section{System Architecture}
\label{sec:architecture}
The application follows a simple layered architecture (Figure~\ref{fig:architecture}). The
user works in a web browser. Every interaction re-runs the Streamlit script \code{app.py},
which shows the tabs for the selected role and calls a function in the corresponding module
file. The module function builds an SQL statement with \code{\%s} placeholders and calls
\code{db.py}, the only file that communicates with PostgreSQL. \code{db.py} uses psycopg to send
the SQL and its parameters to the database server, which checks all constraints, executes the
statement and returns rows or an error. The result travels back up the same path and is
displayed on the screen.

\begin{figure}[H]
  \centering
  \begin{tikzpicture}[
      layer/.style={draw, rounded corners, fill=blue!6, text width=10.5cm,
                    minimum height=1.1cm, align=center, font=\small},
      db/.style={draw, cylinder, shape border rotate=90, aspect=0.25, fill=orange!15,
                 minimum width=7.5cm, minimum height=1.6cm, align=center, font=\small},
      arr/.style={-{Stealth[length=2.5mm]}, thick},
      lbl/.style={font=\scriptsize, fill=white, inner sep=1pt}]
    \node[layer] (browser) {\textbf{Web browser}\\ user clicks, types and selects};
    \node[layer, below=0.9cm of browser] (app)
      {\textbf{app.py} (Streamlit)\\ ``View as'' role selector and tabs};
    \node[layer, below=0.9cm of app] (modules)
      {\textbf{Module files}\\ \code{startup\_module}, \code{application\_module},
       \code{evaluation\_module}, \code{mentor\_module}, \code{progress\_module}};
    \node[layer, below=0.9cm of modules] (dbpy)
      {\textbf{db.py} (psycopg)\\ connect, execute SQL with parameters, commit / roll back,
       translate errors};
    \node[db, below=0.9cm of dbpy] (pg)
      {\textbf{PostgreSQL 18}\\ \code{incubation\_db}: 9 tables, 2 views};

    \draw[arr] ([xshift=-2cm]browser.south) -- node[lbl,left=2pt]{user action} ([xshift=-2cm]app.north);
    \draw[arr] ([xshift=-2cm]app.south) -- node[lbl,left=2pt]{function call} ([xshift=-2cm]modules.north);
    \draw[arr] ([xshift=-2cm]modules.south) -- node[lbl,left=2pt]{SQL + parameters} ([xshift=-2cm]dbpy.north);
    \draw[arr] ([xshift=-2cm]dbpy.south) -- node[lbl,left=2pt]{network (port 5432)} ([xshift=-2cm]pg.north);

    \draw[arr] ([xshift=2cm]pg.north) -- node[lbl,right=2pt]{rows or error} ([xshift=2cm]dbpy.south);
    \draw[arr] ([xshift=2cm]dbpy.north) -- node[lbl,right=2pt]{Python dictionaries} ([xshift=2cm]modules.south);
    \draw[arr] ([xshift=2cm]modules.north) -- node[lbl,right=2pt]{data / messages} ([xshift=2cm]app.south);
    \draw[arr] ([xshift=2cm]app.north) -- node[lbl,right=2pt]{tables, charts} ([xshift=2cm]browser.south);
  \end{tikzpicture}
  \caption{Layered architecture of the application}
  \label{fig:architecture}
\end{figure}

\section{Project Structure}
\label{sec:structure}
Each file has a single responsibility (Table~\ref{tab:structure}).

\begin{table}[H]
  \centering
  \caption{Project files and their responsibilities}
  \label{tab:structure}
  \begin{tabularx}{\textwidth}{@{}lX@{}}
    \toprule
    \textbf{File} & \textbf{Responsibility} \\
    \midrule
    \code{sql/01\_schema.sql}          & Creates the 9 tables, their constraints and the partial unique index \\
    \code{sql/02\_sample\_data.sql}    & Inserts the sample data in a single transaction \\
    \code{sql/03\_constraint\_tests.sql} & Deliberately violates each constraint and demonstrates \code{INSERT}, \code{UPDATE} and \code{DELETE} (rolled back at the end) \\
    \code{sql/04\_queries.sql}         & The report queries Q1--Q19 \\
    \code{sql/05\_views.sql}           & The two views and queries Q20--Q21 that use them \\
    \code{app/db.py}                   & The only file that talks to PostgreSQL \\
    \code{app/app.py}                  & Entry point: role selector and tab layout \\
    \code{app/startup\_module.py}      & Module 1: database functions and screens \\
    \code{app/application\_module.py}  & Module 2, including the admission transaction \\
    \code{app/evaluation\_module.py}   & Module 3 \\
    \code{app/mentor\_module.py}       & Module 4 \\
    \code{app/progress\_module.py}     & Module 5 \\
    \code{tests/test\_app.py}          & Automated test suite (56 test cases) \\
    \bottomrule
  \end{tabularx}
\end{table}
